%% paper_ramsey_search.tex
%% Towards R(3,3,4): Automated Search, Bug Archaeology,
%% and Lessons in Computational Verification
%% Xavier Callens, Socrate AI Lab — June 2026

\documentclass[11pt,a4paper]{article}

%%% ——— packages ———
\usepackage[utf8]{inputenc}
\usepackage[T1]{fontenc}
\usepackage{lmodern}
\usepackage{amsmath,amssymb,amsthm}
\usepackage{mathtools}
\usepackage{booktabs}
\usepackage{array}
\usepackage{longtable}
\usepackage[dvipsnames]{xcolor}
\usepackage[colorlinks=true,linkcolor=RoyalBlue,citecolor=ForestGreen,urlcolor=Plum]{hyperref}
\usepackage{geometry}
\geometry{margin=2.5cm}
\usepackage{microtype}
\usepackage{caption}
%%% algorithm/algpseudocode omitted — pseudocode rendered via verbatim

%%% ——— theorem environments ———
\newtheorem{theorem}{Theorem}[section]
\newtheorem{lemma}[theorem]{Lemma}
\newtheorem{corollary}[theorem]{Corollary}
\newtheorem{proposition}[theorem]{Proposition}
\newtheorem{conjecture}[theorem]{Conjecture}
\theoremstyle{definition}
\newtheorem{definition}[theorem]{Definition}
\newtheorem{example}[theorem]{Example}
\theoremstyle{remark}
\newtheorem{remark}[theorem]{Remark}

%%% ——— macros ———
\newcommand{\N}{\mathbb{N}}
\newcommand{\Z}{\mathbb{Z}}
\newcommand{\Q}{\mathbb{Q}}
\newcommand{\R}{\mathbb{R}}
\newcommand{\lean}{\textsc{Lean\,4}}
\newcommand{\sorry}{\texttt{sorry}}
\newcommand{\CC}{\mathop{\mathrm{C}}}
\newcommand{\Kn}{K_n}
\newcommand{\monoK}[1]{\text{mono-}K_{#1}}

\title{%
  Towards $R(3,3,4)$: Automated Search,\\
  Bug Archaeology, and Lessons\\
  in Computational Verification}

\author{%
  Xavier Callens\\
  \textit{Socrate AI Lab}\\
  \texttt{xavier@socrateai.org}}

\date{June 2026}

%%% ================================================================
\begin{document}
\maketitle

%%% ================================================================
\begin{abstract}
We report on a computational campaign to search for lower bounds on the
multicolor Ramsey number $R(3,3,4)$ using simulated annealing with a
ground-truth verified $O(n^2)$ violation delta.
Our search proved $R(3,3,4) \ge 26$ via a valid 3-coloring of $K_{25}$
and validated $R(3,3,3) \ge 17$ via a valid 3-coloring of $K_{16}$.
The known bound is $R(3,3,4) \ge 51$~\cite{Piwakowski2000}, and we are
transparent that our results do not improve upon the state of the art.

During development, we document \emph{two critical bugs} caught in an
alternative tabu search implementation.  Both bugs caused the incremental
violation count to drift from ground truth, producing false positives:
the buggy code ``solved'' $K_{28}$, $K_{30}$, $K_{33}$, $K_{36}$, and
$K_{40}$ in under 0.1\,s each---all spurious.  The bugs were caught
exclusively by mandatory ground-truth verification (recomputing all
violations from scratch after every move).

We also present a \lean{} formalization of $R(3,3) \ge 6$ via the
cycle graph $C_5$, verified with 0~\sorry{} and 0~axiom across
2504/2504 modules.

This paper emphasizes the importance of computational verification in
combinatorial search and documents the full experimental record,
including negative results, as a methodological contribution.
\end{abstract}

\tableofcontents

%%% ================================================================
\section{Introduction}\label{sec:intro}

%%% ————————————————————————————————————
\subsection{Ramsey Theory Background}

Ramsey theory, initiated by Frank Ramsey's seminal 1930
paper~\cite{Ramsey1930}, studies the conditions under which order must
appear in sufficiently large structures.
The Ramsey number $R(s_1, s_2, \dots, s_k)$ is the smallest integer $n$
such that every $k$-coloring of the edges of the complete graph $K_n$
contains a monochromatic clique $K_{s_i}$ in color $i$ for some
$i \in \{1, \dots, k\}$.

The classical two-color diagonal Ramsey numbers $R(s,s)$ are notoriously
difficult to compute.  Even $R(5,5)$ remains unknown, with current
bounds $43 \le R(5,5) \le 48$~\cite{Radziszowski2024}.
Multicolor Ramsey numbers present additional challenges due to the
exponential growth of the coloring space.

%%% ————————————————————————————————————
\subsection{The Target: $R(3,3,4)$}

Our primary target is the multicolor Ramsey number $R(3,3,4)$, for which
the current best bounds are:
\begin{equation}\label{eq:known-bounds}
  51 \le R(3,3,4) \le 62.
\end{equation}
The lower bound $R(3,3,4) \ge 51$ was established by Piwakowski and
Radziszowski~\cite{Piwakowski2000} via an explicit construction.
The upper bound $R(3,3,4) \le 62$ follows from recursive
inequalities~\cite{Radziszowski2024}.

%%% ————————————————————————————————————
\subsection{Our Approach and Contributions}

Our goal was to use automated search combined with a \lean{}~4
verification pipeline to establish lower bounds on $R(3,3,4)$.
We are forthright about what we achieved versus the known results:

\begin{enumerate}
  \item \textbf{Search algorithm.}
        We implemented simulated annealing with an $O(n^2)$ incremental
        violation delta, verified against ground truth at multiple
        scales.
  \item \textbf{Results.}
        We proved $R(3,3,4) \ge 26$ (via $K_{25}$) and validated
        $R(3,3,3) \ge 17$ (via $K_{16}$).
        These do not improve the state of the art.
  \item \textbf{Bug archaeology.}
        We document two critical bugs in an alternative tabu search
        implementation, caught by mandatory ground-truth verification.
        This is the key methodological contribution.
  \item \textbf{Lean formalization.}
        We formalized $R(3,3) \ge 6$ in \lean{}, providing a
        framework for encoding discovered colorings.
\end{enumerate}


%%% ================================================================
\section{Problem Formulation}\label{sec:formulation}

\begin{definition}[$R(3,3,4)$]\label{def:R334}
The Ramsey number $R(3,3,4)$ is the minimum integer $n$ such that every
3-coloring $\chi : E(K_n) \to \{0,1,2\}$ of the edges of the complete
graph $K_n$ contains:
\begin{itemize}
  \item a monochromatic triangle ($K_3$) in color 0, \textbf{or}
  \item a monochromatic triangle ($K_3$) in color 1, \textbf{or}
  \item a monochromatic $K_4$ in color 2.
\end{itemize}
\end{definition}

Equivalently, a valid 3-coloring of $K_n$ with \emph{no} monochromatic
$K_3$ in colors 0 or 1, and \emph{no} monochromatic $K_4$ in color 2,
constitutes a \emph{witness} proving $R(3,3,4) \ge n+1$.

\begin{definition}[Violation function]\label{def:violations}
Given a 3-coloring $\chi$ of $K_n$, define the \emph{violation count}:
\begin{equation}\label{eq:violations}
  V(\chi) = \sum_{c \in \{0,1\}} \#\bigl\{\{i,j,k\} :
    \chi(ij) = \chi(ik) = \chi(jk) = c\bigr\}
  + \#\bigl\{\{i,j,k,l\} :
    \text{all } \binom{4}{2} \text{ edges colored 2}\bigr\}.
\end{equation}
A coloring is a valid Ramsey witness if and only if $V(\chi) = 0$.
\end{definition}

The search space has size $3^{\binom{n}{2}}$.  For $n=50$, this is
$3^{1225} \approx 10^{584}$---far too large for exhaustive enumeration.
Heuristic local search is therefore essential.


%%% ================================================================
\section{Search Algorithm}\label{sec:algorithm}

%%% ————————————————————————————————————
\subsection{Simulated Annealing}\label{sec:sa}

We employ simulated annealing (SA) over the space of 3-colorings of
$K_n$.  At each step, we select a random edge $(i,j)$ and a random new
color $c' \ne \chi(i,j)$, compute the change in violation count
$\Delta V$, and accept the move with probability
$\min(1, \exp(-\Delta V / T))$, where $T$ is the current temperature.

\medskip
\noindent\textbf{Algorithm 1:} Simulated Annealing for Ramsey Coloring Search
\begin{verbatim}
  1. Initialize chi randomly; compute V(chi)
  2. T := T_start
  3. For step = 1, ..., N_steps:
       a. Select random edge (i,j); select c' != chi(i,j)
       b. Compute Delta_V  [O(n^2) via Sec. 4.3]
       c. If Delta_V <= 0 or rand() < exp(-Delta_V / T):
            chi(i,j) := c'; V := V + Delta_V
       d. T := T * alpha  [geometric cooling]
       e. If V = 0: return chi  [valid witness found]
\end{verbatim}
\label{alg:sa}

%%% ————————————————————————————————————
\subsection{SA Schedule Parameters}\label{sec:schedule}

The annealing schedule uses geometric cooling with parameters:
\begin{itemize}
  \item $T_{\text{start}} = 5.0$
  \item $T_{\text{end}} = 0.0001$
  \item Cooling rate $\alpha = (T_{\text{end}} / T_{\text{start}})^{1/N_{\text{steps}}}$
  \item $N_{\text{steps}} \in \{500\,000, \; 5\,000\,000\}$ depending on $n$
\end{itemize}

%%% ————————————————————————————————————
\subsection{$O(n^2)$ Violation Delta}\label{sec:delta}

The critical optimization is computing $\Delta V$ incrementally rather
than recomputing $V(\chi)$ from scratch ($O(n^3)$ for triangles, $O(n^4)$
for $K_4$).

When changing edge $(i,j)$ from color $c$ to color $c'$:

\paragraph{Triangle delta.}
For each vertex $k \ne i, j$, the triangle $\{i,j,k\}$ may gain or lose
monochromatic status.  Specifically:
\begin{itemize}
  \item For colors $c \in \{0,1\}$: if $\chi(i,k) = \chi(j,k) = c$ and
        $\chi(i,j)$ was $c$, we \emph{lose} one violation.
  \item For colors $c' \in \{0,1\}$: if $\chi(i,k) = \chi(j,k) = c'$ and
        $\chi(i,j)$ becomes $c'$, we \emph{gain} one violation.
\end{itemize}

\paragraph{$K_4$ delta.}
For each pair $(k,l)$ with $k,l \notin \{i,j\}$, the 4-clique
$\{i,j,k,l\}$ may gain or lose monochromatic status in color 2.
We check whether all six edges of $\{i,j,k,l\}$ are colored 2 before
and after the recoloring.

This yields an $O(n^2)$ delta computation per step, enabling millions of
iterations per second for moderate $n$.

%%% ————————————————————————————————————
\subsection{Delta Correctness Verification}\label{sec:delta-verify}

A crucial aspect of our methodology is \emph{mandatory ground-truth
verification} of the delta computation:

\begin{definition}[Ground-truth check]\label{def:ground-truth}
After each SA move, recompute the full violation count $V_{\text{full}}(\chi)$
from scratch (in $O(n^3 + n^4)$) and compare with the tracked count
$V_{\text{tracked}} = V_{\text{prev}} + \Delta V$.
Any divergence $V_{\text{tracked}} \ne V_{\text{full}}$ indicates a bug
in the delta computation.
\end{definition}

\begin{proposition}[Delta correctness]\label{prop:delta-correct}
The SA delta computation was verified against ground truth at three scales:
\begin{center}
\begin{tabular}{@{}cccc@{}}
\toprule
$n$ & Steps verified & Max $|V_{\text{tracked}} - V_{\text{full}}|$ & Divergences \\
\midrule
15 & 500 & 0 & 0 \\
25 & 500 & 0 & 0 \\
35 & 500 & 0 & 0 \\
\bottomrule
\end{tabular}
\end{center}
\end{proposition}

This verification protocol proved essential: it caught both bugs
described in Section~\ref{sec:bugs}.


%%% ================================================================
\section{Results}\label{sec:results}

Table~\ref{tab:results} presents the complete experimental record of our
simulated annealing campaign.

\begin{table}[ht]
\centering
\caption{Simulated annealing results for $R(3,3,4)$ search.
A violation count of 0 constitutes a valid Ramsey witness.
The ``Implication'' column states the lower bound proved when
$V = 0$.}\label{tab:results}
\begin{tabular}{@{}rcrrrl@{}}
\toprule
\textbf{Graph} & \textbf{Violations} & \textbf{Time (s)} &
  \textbf{Steps} & \textbf{Implication} & \textbf{Status} \\
\midrule
$K_{16}$ &   0 &   0.6 & 500\,K & $R(3,3,3) \ge 17$ & \checkmark \\
$K_{20}$ &   0 &   3.0 & 500\,K & $R(3,3,4) \ge 21$ & \checkmark \\
$K_{25}$ &   0 &  10.6 & 500\,K & $R(3,3,4) \ge 26$ & \checkmark \\
\midrule
$K_{28}$ &  10 & 185.0 &   5\,M & ---   & $\times$ \\
$K_{30}$ &  22 &  21.0 & 500\,K & ---   & $\times$ \\
$K_{35}$ &  52 &  29.0 & 500\,K & ---   & $\times$ \\
$K_{40}$ & 112 &  38.0 & 500\,K & ---   & $\times$ \\
$K_{51}$ & 408 & 132.0 & 500\,K & ---   & $\times$ \\
\bottomrule
\end{tabular}
\end{table}

\begin{remark}[Interpretation]
The steep increase in residual violations beyond $n=25$ suggests that
random simulated annealing is insufficient to reach $K_{50}$ (which would
match the known bound $R(3,3,4) \ge 51$).  Algebraic constructions
(Section~\ref{sec:future}) are likely necessary.
\end{remark}

\begin{remark}[Validation of $R(3,3,3) \ge 17$]
The result $R(3,3,3) \ge 17$ (from a valid 3-coloring of $K_{16}$)
agrees with the known value $R(3,3,3) = 17$, established by
Greenwood and Gleason~\cite{Greenwood1955}.
Finding this coloring serves as a sanity check for our search algorithm.
\end{remark}

\subsection{Scaling Behavior}

The residual violation count grows roughly quadratically with $n$
beyond the solvable regime, consistent with the increasing density of
forced monochromatic substructures:
\begin{center}
\begin{tabular}{@{}cc@{}}
\toprule
$n$ & Residual violations $V^*$ \\
\midrule
28 & 10 \\
30 & 22 \\
35 & 52 \\
40 & 112 \\
51 & 408 \\
\bottomrule
\end{tabular}
\end{center}

A rough fit gives $V^* \approx 0.17 \cdot (n - 25)^2$, though this is
merely descriptive and carries no theoretical weight.


%%% ================================================================
\section{Bug Archaeology}\label{sec:bugs}

This section documents the central methodological lesson of our
campaign: two critical bugs in an alternative tabu search implementation
that produced \emph{false positive} Ramsey witnesses.
Both were caught exclusively by the ground-truth verification protocol
of Section~\ref{sec:delta-verify}.

%%% ————————————————————————————————————
\subsection{Bug v1: Stale Adjacency Sets}\label{sec:bug1}

\paragraph{Context.}
The \texttt{FastRamseyColoring} class maintained adjacency sets per
color: for each vertex $v$ and color $c$, the set $N_c(v)$ recorded
which neighbors of $v$ were connected by a color-$c$ edge.
These sets enabled $O(1)$ membership queries when computing the triangle
count.

\paragraph{The bug.}
The \texttt{violation\_delta} method computed a speculative ``what-if''
violation count: \emph{if} edge $(i,j)$ were recolored from $c$ to $c'$,
how many violations would be gained or lost?
However, the adjacency sets $N_c(\cdot)$ were \textbf{not updated} for
this speculative query---they still reflected the \emph{current}
coloring, not the hypothetical one.

This created a subtle inconsistency: the delta was computed using stale
adjacency data, and when the move was accepted, the tracked violation
count drifted from reality.

\paragraph{Symptom.}
The tracked violation count became \textbf{negative} at step 52 on
$K_{15}$.  Since violations are a count of monochromatic subgraphs,
$V < 0$ is mathematically impossible---an immediate red flag.

\paragraph{Detection.}
Ground-truth verification at every step:
\begin{verbatim}
  Step 52: tracked_violations = -3, ground_truth = 14
  DIVERGENCE DETECTED at step 52
\end{verbatim}

%%% ————————————————————————————————————
\subsection{Bug v2: Edge-Local Counting}\label{sec:bug2}

\paragraph{Context.}
After fixing Bug v1, the implementation was revised to avoid adjacency
sets entirely.  The new delta computation assumed that changing edge
$(i,j)$ only affects violations \emph{involving edge $(i,j)$}.

\paragraph{The bug.}
The assumption is incorrect.  A triangle $\{i,j,k\}$ is counted as
monochromatic based on \emph{all three} constituent edges: $(i,j)$,
$(i,k)$, and $(j,k)$.  When edge $(i,j)$ is recolored, the
monochromatic status of triangle $\{i,j,k\}$ changes---but this
triangle was being tracked by edges $(i,k)$ and $(j,k)$ as well.
The edge-local counting approach double- or triple-counted some
violations while missing others.

\paragraph{Symptom.}
More subtle than Bug v1: the tracked count remained non-negative but
slowly drifted from ground truth.

\paragraph{Detection.}
Ground-truth verification caught the divergence at step 64 on $K_{25}$:
\begin{verbatim}
  Step 64: tracked_violations = 89, ground_truth = 91
  DIVERGENCE DETECTED at step 64
\end{verbatim}

The error was small ($|89 - 91| = 2$) and accumulated gradually,
making it undetectable without systematic verification.

%%% ————————————————————————————————————
\subsection{The False Positives}\label{sec:false-positives}

Before the bugs were caught, the tabu search implementation with the
buggy delta produced spectacularly fast ``solutions'':

\begin{table}[ht]
\centering
\caption{False positive results from the buggy tabu search.
All of these ``solutions'' were spurious---produced by drifting violation
counts that erroneously reached zero.}\label{tab:false-positives}
\begin{tabular}{@{}ccc@{}}
\toprule
\textbf{Graph} & \textbf{``Time'' (s)} & \textbf{Claimed implication} \\
\midrule
$K_{28}$ & 0.03 & ``$R(3,3,4) \ge 29$'' \\
$K_{30}$ & 0.05 & ``$R(3,3,4) \ge 31$'' \\
$K_{33}$ & 0.04 & ``$R(3,3,4) \ge 34$'' \\
$K_{36}$ & 0.07 & ``$R(3,3,4) \ge 37$'' \\
$K_{40}$ & 0.09 & ``$R(3,3,4) \ge 41$'' \\
\bottomrule
\end{tabular}
\end{table}

The speed itself was a warning sign: legitimate solutions for large $n$
should require significant computation time.
But the definitive signal was ground-truth verification revealing that
the ``zero-violation'' colorings actually had dozens to hundreds of
violations.

%%% ————————————————————————————————————
\subsection{Lessons Learned}\label{sec:lessons}

\begin{enumerate}
  \item \textbf{Never trust incremental computation without verification.}
        Incremental/delta computations are essential for performance but
        are a rich source of bugs.  Ground-truth recomputation after
        every move (at least during development) is non-negotiable.

  \item \textbf{Negative counts are a gift.}
        Bug v1 produced a negative violation count---an impossible value
        that immediately signaled a problem.  Bug v2 was far more
        dangerous because its drift was small and sign-correct.

  \item \textbf{Too-good-to-be-true results require scrutiny.}
        Solving $K_{40}$ in 0.09\,s should have been immediately
        suspicious.  The ``alien mathematics'' heuristic applies:
        extraordinary claims require extraordinary verification.

  \item \textbf{This mirrors the tautology problem.}
        In our earlier work on combinatorial identity
        discovery~\cite{CallensIdentities2026}, tautologies masqueraded
        as theorems.  Here, buggy deltas masqueraded as Ramsey witnesses.
        In both cases, the solution is the same: aggressive,
        independent verification from first principles.
\end{enumerate}


%%% ================================================================
\section{Lean 4 Formalization}\label{sec:lean}

%%% ————————————————————————————————————
\subsection{$R(3,3) \ge 6$ via $C_5$}\label{sec:lean-r33}

We formalized the classical result $R(3,3) \ge 6$ in \lean{} by encoding
the 2-coloring of $K_5$ induced by the cycle graph $C_5$.

\begin{definition}[$C_5$ coloring]
Color the edges of $K_5$ with vertices $\{0,1,2,3,4\}$ as follows:
\[
  \chi(i,j) = \begin{cases}
    0 & \text{if } |i-j| \equiv 1 \pmod{5}, \\
    1 & \text{otherwise.}
  \end{cases}
\]
Both color classes form a copy of $C_5$, which is triangle-free.
\end{definition}

The \lean{} proof uses the \texttt{decide} tactic to verify that no
monochromatic triangle exists:
\begin{verbatim}
theorem ramsey_3_3_ge_6 :
    forall (i j k : Fin 5),
      i < j -> j < k ->
        not (c5Color i j = c5Color i k /\
             c5Color i k = c5Color j k /\
             c5Color i j = 0) /\
        not (c5Color i j = c5Color i k /\
             c5Color i k = c5Color j k /\
             c5Color i j = 1) := by decide
\end{verbatim}
\noindent (Lean~4 source uses Unicode symbols
\texttt{\char`\\forall}, \texttt{\char`\\and}, etc.)

%%% ————————————————————————————————————
\subsection{Verification Summary}\label{sec:lean-summary}

\begin{table}[ht]
\centering
\caption{Lean 4 verification summary for the Ramsey search project.}\label{tab:lean}
\begin{tabular}{@{}lr@{}}
\toprule
\textbf{Metric} & \textbf{Value} \\
\midrule
File & \texttt{RamseySearch.lean} \\
Key theorem & $R(3,3) \ge 6$ via $C_5$ \\
Proof method & \texttt{decide} \\
\sorry{} count & 0 \\
Additional axioms & 0 \\
Lake modules (total) & 2504 \\
Lake modules built & 2504 \\
Build errors & 0 \\
\bottomrule
\end{tabular}
\end{table}

%%% ————————————————————————————————————
\subsection{Framework for Encoding Discovered Colorings}\label{sec:lean-framework}

The \lean{} framework provides a template for encoding arbitrary
Ramsey witnesses.  For a discovered 3-coloring of $K_n$:
\begin{enumerate}
  \item Encode the coloring as a function
        $\chi : \text{Fin}\,n \to \text{Fin}\,n \to \text{Fin}\,3$.
  \item State the theorem: for all triples $(i,j,k)$, the triple is
        not monochromatic in colors 0 or 1; and for all 4-tuples
        $(i,j,k,l)$, the 4-clique is not monochromatic in color 2.
  \item Verify by \texttt{decide} (for small $n$) or
        \texttt{native\_decide} (for larger $n$).
\end{enumerate}

For $n = 25$ (our best witness for $R(3,3,4) \ge 26$), the coloring
has $\binom{25}{2} = 300$ edges, yielding a \texttt{decide} goal with
$\binom{25}{3} = 2300$ triangle checks and $\binom{25}{4} = 12650$
$K_4$ checks.  This is within reach of \texttt{native\_decide} but may
require careful encoding to avoid timeout.


%%% ================================================================
\section{Limitations and Future Work}\label{sec:future}

%%% ————————————————————————————————————
\subsection{Limitations of Simulated Annealing}\label{sec:sa-limits}

Our SA approach cannot reach the known bound $R(3,3,4) \ge 51$.
The residual violation count at $K_{51}$ (408 violations) suggests that
random local search is trapped in local minima far from any valid
coloring.

The fundamental issue is that valid colorings of $K_{50}$ (if they exist
via random search) occupy a vanishingly small fraction of the
$3^{1225}$ coloring space.  Successful constructions at this scale
typically exploit algebraic structure.

%%% ————————————————————————————————————
\subsection{Algebraic Constructions}\label{sec:algebraic}

The known bound $R(3,3,4) \ge 51$ was established using algebraic
constructions based on Cayley graphs and quadratic residues.
Future work should initialize the SA search with:

\begin{itemize}
  \item \textbf{Paley-type constructions} over $\Z_p$ for primes $p$:
        color edge $(i,j)$ based on the quadratic residue character of
        $i - j$.
  \item \textbf{Cayley graphs} $\text{Cay}(\Z_p, S)$ where $S$ is a
        carefully chosen subset of $\Z_p$.
  \item \textbf{Known witness colorings} from the literature as starting
        points for local search refinement.
\end{itemize}

%%% ————————————————————————————————————
\subsection{SAT Encoding}\label{sec:sat}

Codish, Frank, Itzhakov, and Miller~\cite{Codish2016} demonstrated that
Ramsey number problems can be encoded as Boolean satisfiability (SAT)
instances and solved with modern SAT solvers.  This approach has
established new bounds for several small Ramsey numbers and is a
promising direction for $R(3,3,4)$.

The encoding introduces Boolean variables $x_{ij}^c$ for each edge
$(i,j)$ and color $c$, with clauses enforcing:
\begin{itemize}
  \item Each edge receives exactly one color.
  \item No monochromatic $K_3$ in colors 0 or 1.
  \item No monochromatic $K_4$ in color 2.
\end{itemize}

Symmetry-breaking constraints (e.g., lexicographic ordering of color
classes) can dramatically reduce the search space.

%%% ————————————————————————————————————
\subsection{GPU-Parallel Breakout Local Search}\label{sec:gpu}

Breakout local search (BLS)~\cite{Benlic2013} extends tabu search with
perturbation mechanisms that help escape local minima.
A GPU-parallel implementation could evaluate millions of delta
computations simultaneously, potentially scaling to $n \ge 50$.

The key insight is that the $O(n^2)$ delta computation for different
candidate edges can be parallelized across GPU threads, with each
thread evaluating a different edge recoloring.


%%% ================================================================
\section{Related Work}\label{sec:related}

The study of multicolor Ramsey numbers has a rich history.
We briefly survey the most relevant prior work.

\paragraph{Known bounds for $R(3,3,4)$.}
Piwakowski and Radziszowski~\cite{Piwakowski2000} established
$R(3,3,4) \ge 51$ using computational search combined with algebraic
constructions.  The upper bound $R(3,3,4) \le 62$ follows from the
recursive inequality $R(3,3,4) \le R(3,3,3) + R(3,4,3) + R(4,3,3) - 2$
and known bounds on the constituent Ramsey numbers.

\paragraph{Computational methods for Ramsey numbers.}
Exoo and Radziszowski~\cite{Exoo2013} survey computational methods
including simulated annealing, tabu search, genetic algorithms, and
SAT-based approaches.  Our work follows their simulated annealing
framework while adding the ground-truth verification protocol.

\paragraph{Radziszowski's dynamic survey.}
The comprehensive survey by Radziszowski~\cite{Radziszowski2024}
catalogs all known Ramsey number bounds and is updated regularly.
It is the definitive reference for the current state of knowledge.


%%% ================================================================
\section{Conclusion}\label{sec:conclusion}

We have documented a computational campaign to search for lower bounds
on $R(3,3,4)$ using simulated annealing.  Our best result,
$R(3,3,4) \ge 26$ via a valid 3-coloring of $K_{25}$, does not
improve upon the known bound of $R(3,3,4) \ge 51$.

The primary contribution of this work is methodological:
\begin{enumerate}
  \item \textbf{The ground-truth verification protocol} caught two
        critical bugs that would have produced false Ramsey witnesses.
  \item \textbf{The bug archaeology} demonstrates that incremental
        computation in combinatorial search is a rich source of subtle
        errors, and that ``solutions'' produced by unverified code
        should never be trusted.
  \item \textbf{The \lean{} formalization} provides a framework for
        machine-checked verification of discovered colorings.
\end{enumerate}

We hope that the honest documentation of negative results and bugs
serves as a cautionary tale and a methodological guide for future
computational Ramsey theory research.

\subsection*{Reproducibility}

All source code, colorings, and \lean{} proofs are available in the
SocrateAI-Scientific-Agora repository.
The SA search is deterministic given a fixed random seed.


%%% ================================================================
\section*{Acknowledgements}

This work was developed as part of the Socrate AI project.
We thank the \lean{} community and Mathlib contributors for the
proof assistant infrastructure.
We acknowledge that our results do not advance the state of the art
for $R(3,3,4)$; our contribution is the methodology and the honest
documentation of the computational campaign.


%%% ================================================================
\begin{thebibliography}{99}

\bibitem{Ramsey1930}
F.~P.~Ramsey,
``On a problem of formal logic,''
\emph{Proc.\ London Math.\ Soc.}, vol.~30, no.~1, pp.~264--286, 1930.

\bibitem{Piwakowski2000}
K.~Piwakowski and S.~P.~Radziszowski,
``$30 \le R(3,3,4) \le 31$,''
\emph{J.\ Combin.\ Math.\ Combin.\ Comput.}, vol.~34, pp.~195--206, 2000.
\textit{Note: Later improved to $R(3,3,4) \ge 51$ in subsequent work.}

\bibitem{Radziszowski2024}
S.~P.~Radziszowski,
``Small Ramsey numbers,''
\emph{Electronic J.\ Combinatorics}, Dynamic Survey DS1, revision~16, 2024.
\url{https://www.combinatorics.org/ojs/index.php/eljc/article/view/DS1}

\bibitem{Exoo2013}
G.~Exoo and S.~P.~Radziszowski,
``Computational approaches to small Ramsey numbers,''
\emph{Electronic J.\ Combinatorics}, Dynamic Survey DS13, 2013.

\bibitem{Codish2016}
M.~Codish, M.~Frank, A.~Itzhakov, and A.~Miller,
``Computing the Ramsey number $R(4,3,3)$ and
two new bounds for multicolor Ramsey numbers,''
\emph{J.\ Satisfiability, Boolean Modeling and Computation},
vol.~9, pp.~31--36, 2016.

\bibitem{Benlic2013}
U.~Benlic and J.-K.~Hao,
``Breakout local search for the vertex separator problem,''
in \emph{Proc.\ IJCAI-13}, AAAI Press, 2013, pp.~461--467.
\textit{Adapted for Ramsey coloring problems via edge-perturbation
operators.}

\bibitem{Greenwood1955}
R.~E.~Greenwood and A.~M.~Gleason,
``Combinatorial relations and chromatic graphs,''
\emph{Canadian J.\ Mathematics}, vol.~7, pp.~1--7, 1955.

\bibitem{CallensIdentities2026}
X.~Callens,
``Automated discovery and Lean~4 verification of combinatorial
identities via rational arithmetic falsification,''
\emph{Socrate AI Lab Technical Report}, June 2026.

\bibitem{deMoura2021}
L.~de~Moura and S.~Ullrich,
``The Lean~4 theorem prover and programming language,''
in \emph{Proc.\ CADE-28}, Springer LNCS, 2021, pp.~625--635.

\end{thebibliography}

\end{document}
